\diarytitle{0012}{一次同次線形偏微分方程式の一般解（$y u_x + x u_y = 0$）}

$P(x,y)\,\dfrac{\partial u}{\partial x} + Q(x,y)\,\dfrac{\partial u}{\partial y} = 0$ の一般解は、特性方程式 $\dfrac{dx}{P} = \dfrac{dy}{Q}$ の解 $\varphi(x,y) = C$ を用いて $u = F(\varphi(x,y))$（$F$ は任意関数）と表せる。

特性方程式は
\[
\frac{dx}{y} = \frac{dy}{x}
\quad\Longleftrightarrow\quad
x\,dx = y\,dy
\]
両辺を積分して $\dfrac{x^{2}}{2} = \dfrac{y^{2}}{2} + C_{1}$、すなわち
\[
x^{2} - y^{2} = C
\]
が特性曲線である。したがって一般解は、任意関数 $F$ を用いて
\[
\boxed{\; u(x,y) = F(x^{2} - y^{2}) \;}
\]
（検算: $u_x = 2x\,F'$, $u_y = -2y\,F'$ より $y u_x + x u_y = 2xy\,F' - 2xy\,F' = 0$。）
